\documentclass[10pt,a4paper]{amsart}
\usepackage[margin=19mm]{geometry}
\usepackage{amsmath,amssymb,mathtools}
\usepackage[scr=rsfs,cal=euler]{mathalfa}
\usepackage{xfrac}
\usepackage[unicode,hidelinks,bookmarks=false]{hyperref}
\newcommand{\ie}{i.e.\ }
\newcommand{\cf}{cf.\ }
\newcommand{\resp}{respectively\ }
\newcommand{\Subsection}{\subsection}
\providecommand{\nobreakdash}{\nobreak\hbox{-}}
\setlength{\emergencystretch}{2em}
\multlinegap0pt
\usepackage{bm}
\usepackage{bbm}
\usepackage{dsfont}
\usepackage{xspace}
\usepackage{cancel}
\usepackage{mathtools}
\usepackage{amsthm}
\makeatletter
\@ifundefined{lemma}{\newtheorem{lemma}{Lemma}}{}
\@ifundefined{observation}{\newtheorem{observation}{Observation}}{}
\makeatother
\DeclareMathOperator{\per}{per}
\newcommand{\del}{\backslash}
\title[An Algorithmic Upper Bound for Permanents via a Permanental ]{An Algorithmic Upper Bound for Permanents via a Permanental Schur Inequality}
\author{Aditi Laddha, Madhusudhan Reddy Pittu}
\date{Source: arXiv:2509.08121v1 (CC BY 4.0)}
\begin{document}
\maketitle

\noindent\emph{Source and licence.} An Algorithmic Upper Bound for Permanents via a Permanental Schur Inequality. Version arXiv:2509.08121v1 (CC BY 4.0), distributed under the Creative Commons Attribution 4.0 International licence, as recorded in the arXiv licence metadata for this exact version and shown on the abstract page https://arxiv.org/abs/2509.08121v1. Full rights notice: licenses/arxiv-2509.08121-CC-BY-4.0.txt.\par\smallskip

\noindent\textit{Editorial scope.} Counted unit: the Lemma whose source label is \texttt{lem:cycles\_bound} in the article (source0.tex lines 1144--1151, source label \texttt{lem:cycles\_bound}), delivered together with the article's own complete original proof of it (source0.tex lines 1153--1281, one \texttt{proof} environment of 129 lines). No mathematics is written, repaired or re-derived: statement and proof are the article's own text line for line. Required context retained uncounted: obs:perm\_ratio\_bound (lines 1130--1136). Every definition, display and label of those lines is retained unchanged. Omitted from this package and not redistributed: 1--1129, 1137--1143, 1152--1152, 1282--1401. Nothing inside a retained excerpt is omitted or altered: every retained body is delivered exactly as the article's lines are written, so no omission receipt of any kind accompanies this package. Citations: the retained excerpts carry no citation command at all, so no bibliography entry of the article is transcribed and no reference is redistributed. Reference adaptation: none was needed and none was made; every reference inside the delivered unit resolves inside it, so no unresolved reference or citation is produced. Printed slips kept literal: no printed word, symbol, hypothesis or number of the counted statement or of its proof was altered. Layout and class: the article's own class is \texttt{article} with packages including fontenc, babel, geometry, setspace, footmisc, framed, parskip, amsmath, amsthm, amssymb, mathtools, bm, bbm, nicefrac; the wrapper is the lane's pinned \texttt{amsart} editorial wrapper at 10pt on A4, which restates the theorem-like environments the retained text needs and adds the article's own macro definitions that the retained text needs (\texttt{\textbackslash del}, \texttt{\textbackslash per}), with extra wrapper packages (bm, bbm, dsfont, xspace, cancel, mathtools). The delivered numbering is the unit's own. Assets: no retained span contains a figure environment, a graphics-inclusion command, a PDF-page inclusion, a table or a TikZ picture; no image file is copied into this package, the package declares no asset dependency and no image file is redistributed with it. Licence: version arXiv:2509.08121v1 of this article is distributed under the Creative Commons Attribution 4.0 International licence, as recorded in the arXiv licence metadata for this exact version and shown on the abstract page https://arxiv.org/abs/2509.08121v1. A full rights notice is delivered at licenses/arxiv-2509.08121-CC-BY-4.0.txt. Characters: every retained excerpt is pure ASCII, so the delivered engine, XeLaTeX with its default Latin Modern text font, reports no missing character.
\begin{observation} 
\label{obs:perm_ratio_bound}
Let $A$ be an $n \times n$ matrix with $a_{i,i} = 1$ for all $i \in [n]$ and $0\leq a_{i,j} \le M$ for all $i$ for some $1\leq M$. Then for every $ S\subseteq [n]$ and $i,j \in [n]\del S$
\begin{equation*}
    \frac{\per(A(S+i, S+j))}{\per(A(S,S))} \le \gamma_{|S|+1}.
\end{equation*}
\end{observation} 

\begin{lemma}\label{lem:cycles_bound}
Let $A$ be an $n \times n$ matrix with $a_{i,i} = 1$ for all $i \in [n]$ and $0\leq a_{i,j} \le M$ for all $i,j$ for some $1\leq M$. Let $A^{(t)}$ be the matrix after $t-1$ steps of the permanent process. For any set of indices $S \subseteq \{t+1, \dots, n\}$ and a fixed element $i_0 \in S$, let $\mathcal{C}$ be the set of all cyclic permutations of elements in $S$. Then
\begin{equation*}
     \frac{\displaystyle \sum_{\sigma \in \mathcal{C}} \prod_{i \in S} a^{(t)}_{i,\sigma(i)}}
       {\per\!\big(A^{(t)}(S - i_0 ,\,S - i_0)\big)}
  \;\leq\; B(n,|S|,t).
\end{equation*}
\end{lemma}

\begin{proof}
    We will prove this by induction on $t$. For the base case, at $t = 1$, $A^{(1)} = A$ and by Observation \ref{obs:perm_ratio_bound},
    \begin{equation*}
          \frac{\sum_{\sigma \in \mathcal{C}} \prod_{i \in S} a_{i, \sigma(i)}}{\per(A_{S\backslash \{i_0\}, S \backslash \{i_0\}})} \leq \frac{\per(A_{S, S})}{\per(A_{S\backslash \{i_0\}, S \backslash \{i_0\}})} \leq |S|! \cdot M^{|S|} = \gamma_{|S|}\leq B(n,|S|,1).
    \end{equation*}

Without loss of generality, let $S=\{t+1, \ldots, t+k\}$ and let $i_0 = t+k$. 
    For $t > 1$, we have $a^{(t)}_{i,j} = a^{(t-1)}_{i,j} + \frac{a^{(t-1)}_{i,t} \cdot a^{(t-1)}_{t,j}}{a^{(t-1)}_{t,t}}$. 
    For ease of notation, we will use $b_{i,j}$ to denote $a^{(t-1)}_{i,j}$.  Then, we have

    \begin{align*}
        \frac{\sum_{\sigma \in \mathcal{C}} \prod_{i = t+1}^{t+k} a_{i, \sigma(i)}}{\per(A_{S\backslash \{t+k\}, S \backslash \{t+k\}})} &=  \frac{\sum_{\sigma \in \mathcal{C}} \prod_{i = t+1}^{t+k} \left(b_{i, \sigma(i)} + \frac{b_{i,t} b_{t,\sigma(i)}}{b_{t,t}}\right)}{\per\left(B_{S\backslash \{t+k\}, S \backslash \{t+k\}} + \frac{b_{S\backslash \{t+k\}, t} b_{t, S\backslash \{t+k\}}}{b_{t,t}}\right)} \\
        &= \frac{\sum_{\sigma \in \mathcal{C}} \prod_{i = t+1}^{t+k} \left(b_{i, \sigma(i)} b_{t,t} + b_{i,t} b_{t,\sigma(i)}\right)}{b_{t,t} \per\left(b_{t,t} \cdot B_{S\backslash \{t+k\}, S \backslash \{t+k\}} +b_{S\backslash \{t+k\}, t} b_{t, S\backslash \{t+k\}}\right)}.
    \end{align*}
    We define the numerator as a polynomial in $b_{t,t}$ as follows:
    \begin{equation*}
        \sum_{\sigma \in \mathcal{C}} \prod_{i = t+1}^{t+k} \left(b_{i, \sigma(i)} b_{t,t} + b_{i,t} b_{\sigma(i), t}\right) := \sum_{\ell = 0}^{k} \alpha_{\ell} \cdot b_{t,t}^{k-\ell}.
    \end{equation*}
    Without loss of generality, assume that $k$ is even. The proof for odd $k$ follows similarly. Let $D$ denote the denominator, we will show that 
    \begin{align*}
        \sum_{\ell \in \{0, 2, \ldots, k-2, k\}} \alpha_{\ell} \cdot b_{t,t}^{k-\ell} &\leq B(n, k, t-1) \ \cdot D \\
        \sum_{\ell \in \{1, 3, \ldots, k-3, k-1\} } \alpha_{\ell} \cdot b_{t,t}^{k-\ell} &\leq B(n, k+1, t-1) \cdot D.
    \end{align*}
    Combining the two equations gives
    \begin{align*}
        \sum_{\ell = 0}^{k} \alpha_{\ell} \cdot b_{t,t}^{k-\ell} \leq \left( B(n, k, t-1) + B(n, k+1, t-1)\right) \cdot D = B(n, k, t)  \cdot D,
    \end{align*}
where we used $B(n, a, b-1) + B(n, a+1, b-1) = B(n, a, b)$.

We will also expand the denominator $D$ also as a polynomial in $b_{t,t}$ as follows:
 \begin{align*}
        D = b_{t,t} \cdot \per\left(b_{t,t} \cdot B_{S\backslash \{t+k\}, S \backslash \{t+k\}} +b_{S\backslash \{t+k\}, t} b_{t, S\backslash \{t+k\}} \right) := \sum_{\ell = 0}^k \beta_\ell \cdot b_{t,t}^{k-\ell +1}.
    \end{align*}

    We will now define $\beta_{\ell}$ using the notion of cycle covers in a certain graph. 
    Consider a directed complete graph on the vertex set $S \cup \{t\}$. For $\ell \in \{0, \ldots, k\}$ and $X \subseteq S$, let $\mathcal{CC}_{\geq \ell}(X)$ denote the set of cycle covers of $X \cup \{t\}$ in this graph satisfying the following properties:
    \begin{itemize}
        \item every vertex in $X$  has exactly one incoming and one outgoing edge,
        \item there are exactly $\ell$ cycles containing $t$, and
        \item $t$ does not have a self-loop.
    \end{itemize}
    
    Then
    \begin{equation*}
        \beta_\ell = \ell! \cdot \sum_{C \in \mathcal{CC}_{\geq \ell}(S \backslash \{t+k\})} \prod_{(i\rightarrow j) \in C} b_{i,j}.
    \end{equation*}
    Here, $\ell!$ accounts for the fact that given cover $C \in \mathcal{CC}_{\geq \ell}(S)$, every matching of the $\ell$ incoming edges at $t$ to the $\ell$ outgoing edges at $t$ in $C$ gives a valid permutation on $S$.
    
    We will now relate $\alpha_\ell$ with $\beta_\ell$ using a different notion of cycle cover. For a fixed $\ell \in \{0, \ldots, k\}$ and $X \subseteq S$, let $\mathcal{CC}_{=\ell}(X)$ denote the set of all cycle covers of $X \cup \{t\} $ satisfying the following properties:
    \begin{itemize}
        \item there are no self-loops,
        \item every vertex in $X$ has exactly one incoming and one outgoing edge, and
        \item there are exactly $\ell$ cycles and each cycle contains $t$.
    \end{itemize}
    
    
  Then $ \alpha_{\ell}$ can be stated as
    \begin{align}
        \alpha_{\ell} = (\ell-1)!\sum_{C\in \mathcal{CC}_{=\ell}(S)} \prod_{(i\rightarrow j) \in C} b_{i,j} . \label{eq:alpha}
    \end{align}
    Here, we have an $(\ell-1)!$ term because the number of cyclic permutations in $\mathcal{C}$ which can lead to a cover $C \in \mathcal{CC}_{=\ell}$ is equal to the number of ways to arrange $\ell$ elements in a cycle, since $t$ has $\ell$ incoming and outgoing edges in $C$.

    Observe that there is a unique way to decompose any $C \in \mathcal{CC}_{\geq \ell}(S)$ or $C \in \mathcal{CC}_{=\ell}(S)$ into edge-disjoint cycles. So, for a cycle cover $C \in \mathcal{CC}_{=\ell}(S)$, let $X$ denote the vertex set of the cycle in $C$ containing $t+k$. We can re-write the equation \eqref{eq:alpha} in terms of $X$ as follows.
   \begin{align}
        \alpha_{\ell} &= (\ell-1)! \sum_{X \subseteq S: t+k \in X} \left(\sum_{c \text{ cycle on } X} \prod_{(i\rightarrow j) \in c} b_{i,j}\right) \cdot \left(   \sum_{C' \in \mathcal{CC}_{=\ell-1}(S \backslash X)} \prod_{(i\rightarrow j) \in C'} b_{i,j} \right). \label{eq:recursion}
    \end{align}

By the inductive hypothesis, we have 
\begin{equation*}
    \left(\sum_{c \text{ cycle on } X} \prod_{(i\rightarrow j) \in c} b_{i,j}\right) \leq B(n, |X|, t-1) \cdot \per(B_{X \backslash \{t+k\}, X \backslash \{t+k\}}).
\end{equation*}
Substituting this in equation \eqref{eq:recursion} gives
    \begin{align*}
        \alpha_{\ell} 
        &\leq (\ell-1)! \sum_{X \subseteq S: t+k \in X} B(n, |X|, t-1) \cdot \per(B_{X \backslash \{t+k\}, X \backslash \{t+k\}})  \cdot \left(   \sum_{C' \in \mathcal{CC}_{=\ell-1}(S \backslash X)} \prod_{(i\rightarrow j) \in C'} b_{i,j} \right).
    \end{align*}
For $\ell \neq 1$, $X$ can contain at most $k$ vertices and for $\ell = 1$, $|X| = k+1$, so for any even $\ell$, we have 
 \begin{align*}
        \alpha_{\ell} 
        &\leq (\ell-1)! \cdot B(n, k, t-1)  \sum_{X \subseteq S: t+k \in X}  \per(B_{X \backslash \{t+k\}, X \backslash \{t+k\}})  \cdot \left(   \sum_{C' \in \mathcal{CC}_{\ell-1}(S \backslash X)} \prod_{(i\rightarrow j) \in C'} b_{i,j} \right).
    \end{align*}

For a set $X$, 
\begin{align*}
    \per(B_{X \backslash \{t+k\}, X \backslash \{t+k\}}) &= \sum_{\sigma \in \mathcal{S}_{|X|-1}} \prod_{i \in X\backslash \{t+k\}} b_{i, \sigma(i)} \\
    &= b_{t,t} \sum_{C \in \mathcal{CC}_{\geq 0}(X \backslash \{t+k\})} \prod_{(i\rightarrow j) \in C} b_{i,j}+ \sum_{C \in \mathcal{CC}_{\geq 1}(X\backslash \{t+k\})} \prod_{(i\rightarrow j) \in C} b_{i,j}.
\end{align*}
The second equation follows by dividing the permutations on $X$ based on whether they contain the self-loop $t \rightarrow t$.

Therefore, 
\begin{align}
        \alpha_{\ell} 
        &\leq (\ell-1)! \cdot B(n, k, t-1)  \sum_{X \subseteq S: t+k \in X}     \left( b_{t,t} \sum_{C \in \mathcal{CC}_{\geq 0}(X \backslash \{t+k\})} \prod_{(i\rightarrow j) \in C} b_{i,j}\right)  \sum_{C' \in \mathcal{CC}_{=\ell-1}(S \backslash X)} \prod_{(i\rightarrow j) \in C'} b_{i,j} \notag\\
        &+(\ell-1)! \cdot B(n, k, t-1)  \sum_{X \subseteq S: t+k \in X}    \left( \sum_{C \in \mathcal{CC}_{\geq 1}(X\backslash \{t+k\})} \prod_{(i\rightarrow j) \in C} b_{i,j}\right)  \sum_{C' \in \mathcal{CC}_{=\ell-1}(S \backslash X)} \prod_{(i\rightarrow j) \in C'} b_{i,j} . \label{eq:big}
    \end{align}
    We will bound the two summations in equation \eqref{eq:big} separately.
    
    
    Note that for any $C\in \mathcal{CC}_{\geq 0}(X \backslash \{t+k\})$ and $C' \in \mathcal{CC}_{= \ell-1}(S\backslash X)$, given $C \cup C'$, there is only one way to reconstruct the set $X$: $X$ must contain $t$ and all vertices in cycles not containing $t$, i.e., vertices in $C$. Therefore, 
    \begin{align*}
         \sum_{X \subseteq S: t+k \in X}   &\left( b_{t,t} \sum_{C \in \mathcal{CC}_{\geq 0}(X \backslash \{ t+k\})} \prod_{(i\rightarrow j) \in C} b_{i,j}\right)  \sum_{C' \in \mathcal{CC}_{=\ell-1}(S \backslash X)} \prod_{(i\rightarrow j) \in C'} b_{i,j} \\
         &= b_{t,t} \sum_{C\in \mathcal{CC}_{\geq \ell-1}(S \backslash \{t+k\}) } \prod_{(i\rightarrow j) \in C} b_{i,j} = b_{t,t} \beta_{\ell-1}.
    \end{align*}
    However, given $C \cup C'$ with $C\in \mathcal{CC}_{\geq 1}(X \backslash \{t+k\})$ and $C' \in \mathcal{CC}_{= \ell-1}(S\backslash X)$, there are at most $\ell$ possibilities for set $X$: $X$ must contain all vertices in any cycle not passing through $t$, and $X$ must contain vertices from exactly one cycle passing through $t$. Since $C \cup C'$ contains $\ell$ cycles passing through $t$, we have
    \begin{align*}
         \sum_{X \subseteq S: t+k \in X}  \sum_{C \in \mathcal{CC}_{\geq 1}(X \backslash \{t+k\})} \prod_{(i\rightarrow j) \in C}& b_{i,j} \cdot  \sum_{C' \in \mathcal{CC}_{\ell-1}(S \backslash X)} \prod_{(i\rightarrow j) \in C'} b_{i,j} \\
         &\leq \ell\cdot  \sum_{C\in \mathcal{CC}_{\geq \ell}(S \backslash \{t+k\}) } \prod_{(i\rightarrow j) \in C} b_{i,j} 
         = \ell \cdot \beta_{\ell}.
    \end{align*}
    Plugging these bounds in equation \eqref{eq:big}, we get 
    \begin{align*}
        \alpha_{\ell} 
        &\leq B(n, k, t-1) \cdot  \left( b_{t,t} \beta_{\ell-1} + \beta_{\ell}\right).
    \end{align*}
    
    
    Summing over all even $\ell$'s, we have
    \begin{align}
        \sum_{\ell \in \{0, 2, \ldots, k-2, k\}} \alpha_{\ell} \cdot b_{t,t}^{k-\ell} &\leq B(n, k, t-1) \cdot \left(\beta_0 \cdot b_{t,t}^k + \sum_{\ell \in \{2, \ldots, k-2, k\}} \beta_{\ell-1} b_{t,t}^{k-\ell+1} +   \beta_{\ell} \cdot b_{t,t}^{k-\ell} \right) \notag\\
        &= B(n, k, t-1)  \cdot D. \label{eq:bounded-1}
    \end{align}

    Similarly, for odd $\ell$'s, we have
 \begin{align}
        \sum_{\ell \in \{1, 3, \ldots, k-1\}} \alpha_{\ell} \cdot b_{t,t}^{k-\ell} &\leq B(n, k+1, t-1) \cdot \left(\sum_{\ell \in \{1, 3, \ldots, k-1\}} \beta_{\ell-1} b_{t,t}^{k-\ell+1} +   \beta_{\ell} \cdot b_{t,t}^{k-\ell} \right) \notag\\
        &= B(n, k+1, t-1) \cdot D.\label{eq:bounded-2}
    \end{align}
    Summing equation \eqref{eq:bounded-1} and equation \eqref{eq:bounded-2} completes the proof.
\end{proof}

\end{document}
